% -----------------------------------------------------------------------------
% This file is automatically generated.
% Do not edit manually.
% Source: notebooks/support/hypothesis-testing/support.Rmd
% -----------------------------------------------------------------------------


\noindent To build a margin for one error, we may increase the critical region to $\{k | k \leq 1\}$, resulting in a sample size of $n =$ 473.

\par\addvspace{\topsep}
\begingroup
\fvset{listparameters={%
  \setlength{\topsep}{0pt}%
  \setlength{\partopsep}{0pt}%
  \setlength{\parsep}{0pt}%
  \setlength{\itemsep}{0pt}%
}}
\begin{Verbatim}[breaklines=true,formatcom=\color{ada_blue}]
myAttSample.size(c=1).n
\end{Verbatim}
\begin{Verbatim}[breaklines=true,formatcom=\color{ada_light_blue}]
[1] 473
\end{Verbatim}
\endgroup
\par\addvspace{\topsep}
\noindent In MUS, we increase the critical region by anticipating on the expected error (in monetary terms) in the population. Thus, in applications where selected items are either completely correct or completely incorrect, it is easier to calculate the minimum required sample size using the \ttblue{att\_obj} than using the \ttblue{mus\_obj}.
